% paper/sections/07_time_integration.tex
% §7: 時間積分：物理項ごとの剛性分離と合成
% ═══════════════════════════════════════════════════════════════════════════════
\section{時間積分スキーム：物理項ごとの剛性分離と合成}
\label{ch:time_integration}
\label{sec:time_int}
% ═══════════════════════════════════════════════════════════════════════════════

\noindent\textbf{本章の位置付け：}
本章は，第~\ref{sec:onefluid}章で導出した一流体 NS 方程式
（移流・粘性・圧力・界面張力・浮力の各項；式~\eqref{eq:NS_onefluid}）と
第~\ref{sec:cls_advection}節で確立した CLS 保存形移流方程式を，
第~\ref{sec:CCD}章の CCD/FCCD/UCCD6 空間離散化と整合する形で
1 物理時間ステップに閉じるための時間離散化を与える．
設計の中心は，すべての項を一つの万能積分子へ押し込むことではなく，
各物理項が生む固有の時間スケールを分離し，剛性項だけを陰的に扱い，
非剛性項を陽的に外挿することである．

1 ステップの情報の流れは次の順に固定する．まず直前ステップで投影済みの速度
$\bu^n$ を用いて CLS を更新し，$\psi^{n+1}$ から
$\rho^{n+1}$ と $\kappa_{lg}^{n+1}$ を確定する．次に NS 予測子で
対流・粘性・浮力を処理し，最後に IPC 形式の PPE で非圧縮拘束と
界面張力圧力ジャンプを同時に閉じる．この順序により，界面位置，
物性値，予測速度，圧力増分が同じ物理時刻で接続される．

本稿が採用する項別配分は次表である：

\begin{tcolorbox}[colback=blue!3!white,colframe=blue!50!black,
                  title=\textbf{本稿の項別時間積分},
                  fonttitle=\small\bfseries,boxrule=0.4pt,arc=2pt]
\small
\renewcommand{\arraystretch}{\PaperTableArrayStretch}
\begin{tabular}{@{}p{0.30\linewidth}p{0.40\linewidth}p{0.22\linewidth}@{}}
\hline
物理項 & 時間積分 & 本章の節 \\
\hline
界面（CLS）移流 & TVD-RK3（SSPRK3） & §\ref{sec:time_level1} \\
NS 対流 & IMEX-BDF2（EXT2 外挿） & §\ref{sec:time_level1} \\
粘性 & implicit-BDF2 + 欠陥補正（DC） & §\ref{sec:time_viscous} \\
界面張力 & 表面エネルギー仮想仕事に基づく毛管閉包 + 成分体積制約付き鞍点系 & §\ref{sec:time_csf} \\
圧力投影 & IPC + FCCD PPE & §\ref{sec:ipc_derivation} \\
浮力（上昇気泡） & Balanced--Force 浮力分解（陽的残差） & §\ref{sec:time_buoyancy} \\
\hline
\end{tabular}
\end{tcolorbox}

表は\textbf{物理項ごとの時間処理}を分類したものであり，
1 ステップ内の因果順序そのものではない．実際の順序は前段で述べた通り
CLS 保存形更新 $\to$ NS 予測子 $\to$ IPC PPE 投影である．
一方，以下の説明順では，表面張力を圧力ジャンプとして PPE に内蔵する構造を先に示し，
それを受けて IPC PPE 投影を定義した後，浮力の静水圧成分が
同じ PPE の楕円型作用素に吸収されることを説明する．
このため，章内の説明順は「因果順序」ではなく「閉包依存関係」を反映している．

\begin{tcolorbox}[defbox, title={保存形共通フラックス経路の時間順序}]
保存形共通フラックスを選ぶ場合は，表の「CLS 移流」と「NS 対流」は
独立した二つの陽的対流処理ではない．
TVD-RK3 の各段階で $q$ の面フラックスを作り，
同じ段階台帳から $m$ と $\bm{p}_m$ を更新する：
\[
  (q,m,\bm{p}_m)^n
  \xrightarrow{\mathrm{TVD\text{-}RK3\ flux}}
  (q,m,\bm{p}_m)^T .
\]
その後の NS 予測子は，輸送済みの
$\bu^T=\bm{p}_m^T/m^T$ に対して粘性，重力残差，圧力履歴，
毛管補正などの非移流項を作用させる．
この経路では速度主変数対流の AB2/EXT2 履歴を使わず，
速度履歴は BDF2 の時間導関数に必要な代表値としてのみ用いる．
そうしないと，保存形の
$\partial_t\bm{p}_m+\bnabla\cdot(\bm{p}_m\otimes\bu)$ と，
速度主変数形式の $\bu\cdot\bnabla\bu$ が同じ非線形仕事を二重に与える．
\end{tcolorbox}

\noindent\textbf{時間積分の 3 形式：}
本章で扱う時間積分は，未知ベクトル $u^{n+1}$ の取り扱いにより
\textbf{陽解法（explicit）}・\textbf{陰解法（implicit）}・
\textbf{IMEX 法（implicit-explicit）}の 3 形式に分類される．
陽解法は $u^{n+1}$ を $u^n$ の関数のみで定義するため代数方程式を解かずに前進できるが，
半離散演算子の固有値に対し時間刻みの絶対安定領域が制限される．
陰解法は $u^{n+1}$ を未知側に置く線形（または非線形）求解を要するが，
A-stable 系であれば左半平面の固有値に対して任意 $\Delta t$ で絶対安定となる．IMEX 法は剛性項のみを
陰側へ配し，残余を陽的外挿で前進させる混合形式である．
本章では，この分類を
\textbf{時間スケールの識別 → 項別の積分形式 → 設計次数と実効次数の分離}
の順に用いる．

\noindent\textbf{本章の構成：}
本章の組織化原則は，各物理項を独立な剛性源と見なし，
項ごとに時間処理を選択することにある．具体的には：

\begin{itemize}[leftmargin=2em]
  \item \textbf{移流項}：CLS 移流に TVD-RK3，NS 対流に IMEX-BDF2（EXT2 外挿）を用い，
        移流由来の陽的制約を対流 CFL として残す（§\ref{sec:time_level1}）．
  \item \textbf{粘性項}：implicit-BDF2 + 欠陥補正（DC）で全応力 Helmholtz 系を解き，
        粘性 CFL を時間刻み制約から外す（§\ref{sec:time_viscous}）．
  \item \textbf{界面張力項}：$J_p=j_{gl}=-\sigma\kappa_{lg}$ は
        Young--Laplace 代表として PPE の切断面圧力ジャンプ補正に置く．
        標準毛管力は再初期化前の移流後状態における表面エネルギーの仮想仕事から
        面上の補正量として構成し，圧力仕事と随伴な成分体積制約付き鞍点系で
        静的反力と非平衡駆動を分離してから圧力補正と同じ時刻で結合する
        （§\ref{sec:time_csf}）．
  \item \textbf{圧力投影}：IPC により $\delta p = p^{n+1} - p^n$ を更新し，
        FCCD PPE で非圧縮拘束を同一面上で閉じる（§\ref{sec:ipc_derivation}）．
  \item \textbf{浮力項}：静水圧成分を PPE の楕円型作用素に吸収し，
        残差成分だけを陽的に加える（§\ref{sec:time_buoyancy}）．
\end{itemize}

本章の節構成は：
時間積分形式の分類と精度整合性（§\ref{sec:time_accuracy_principle}），
演算子別剛性マップ + CFL 理論体系（§\ref{sec:time_operator_stiffness}），
移流（§\ref{sec:time_level1}）→粘性（§\ref{sec:time_viscous}）
→界面張力（§\ref{sec:time_csf}）→圧力投影（§\ref{sec:ipc_derivation}）
→浮力（§\ref{sec:time_buoyancy}），
設計次数と検証上の実効次数（§\ref{sec:time_accuracy_table}），
時間刻み制御（§\ref{sec:time_guide}）の順である．

% -------------------------------------------------------------------------------
\subsection{時間積分形式の分類と精度整合性}
\label{sec:time_accuracy_principle}
% -------------------------------------------------------------------------------

\subsubsection*{時間積分形式の分類（explicit / implicit / IMEX）}
\label{sec:scheme_classification}

時間方向 1 ステップ更新を $u^{n+1} = \Phi_{\Delta t}(u^n, u^{n-1}, \ldots; u^{n+1})$
と書くとき，未知 $u^{n+1}$ の取扱いに応じて以下の 3 分類が得られる：
\begin{itemize}[leftmargin=2em]
  \item \textbf{陽解法（explicit）}：
        $u^{n+1} = u^n + \Delta t\,f(u^n)$．未知が右辺に現れず
        ステップごとの非線形求解を要しない．時間刻みは絶対安定領域
        $\{z\in\mathbb{C} : |R(z)|\le 1\}$（$R$ は安定関数）で制約され，
        対流項なら $|\bu|\Delta t/h\lesssim C_{\mathrm{adv}}$ の CFL を生む．
  \item \textbf{陰解法（implicit）}：
        $u^{n+1} = u^n + \Delta t\,f(u^{n+1})$．各ステップで線形・非線形求解
        が必要となるが，BDF1/BDF2 などの A-stable 系は
        任意 $\Delta t$ で絶対安定．
        粘性 CFL $\nu\Delta t/h^2$ などの剛性 CFL 制約を消滅させうる．
  \item \textbf{IMEX 法（implicit--explicit）}：
        $u^{n+1} - u^n = \Delta t\,[\,f_E(u^n,\ldots) + f_I(u^{n+1})\,]$．
        剛性項 $f_I$ のみを陰側に置き，非剛性 $f_E$ を陽的外挿で前進させる
        混合形式（IMEX 一般論は~\cite{AscherRuuthSpiteri1997,AscherRuuthWetton1995}）．本稿の主軸 IMEX-BDF2（EXT2 外挿）はこの典型である．
\end{itemize}
本稿は，CLS 移流が陽的 TVD-RK3 / NS 対流が陽的 EXT2 外挿 /
粘性が陰的 BDF2 / 圧力が投影陰解法（IPC + FCCD PPE）/ 表面張力が
PPE 内蔵の圧力ジャンプ補正量 / 浮力が陽的残差評価の混合構成である．
各項の選択根拠は §\ref{sec:time_operator_stiffness}（剛性マップと CFL 数）に，
零安定多段法の精度整合性は本節後段に与える．

\subsubsection*{精度整合性と零安定性（Lax 等価定理 + Dahlquist 根条件）}
\label{sec:time_accuracy_lax}

本章の目的は，CCD/FCCD/UCCD6 が与える高次空間離散を時間方向で
不整合に壊さないことである．二相流の一流体式では，圧力投影，
界面位置更新，物性値更新，表面張力ジャンプが 1 ステップ内で連鎖するため，
時間次数は単独の常微分方程式時間離散の次数だけでは決まらない．
本稿ではこの連鎖を
\[
  \psi^{n+1}\to (\rho^{n+1},\kappa_{lg}^{n+1})
  \to \bu^* \to (\delta p,\bu^{n+1},p^{n+1})
\]
の同時刻情報流として閉じ，均質流域で $\Ord{\Delta t^2}$，
界面移流単独で $\Ord{\Delta t^3}$ の時間精度を目標とする．
高次空間離散は空間誤差と寄生流を抑える役割を担い，時間方向の律速は
表~\ref{tab:time_accuracy_prod} の項別合成で明示する．

\noindent\textbf{Lax 等価定理と零安定性：}
線形定数係数の well-posed な初期値問題に対し，
\textbf{一貫性}（local truncation error $\to 0$ as $\Delta t \to 0$）と
\textbf{零安定性}（zero-stability; Dahlquist の根条件）の同時成立が
\textbf{収束}（$\Ord{\Delta t^q}$）の必要十分条件である~\cite{HairerWanner1996}．
多段法 $\sum_{k=0}^{s}\alpha_k u^{n+k} = \Delta t \sum_{k=0}^{s}\beta_k f(u^{n+k})$
の特性多項式 $\rho(\zeta) = \sum_{k=0}^{s}\alpha_k \zeta^k$ について，
全根が単位円内 $|\zeta|\le 1$ かつ単位円上の根は単純，が零安定の必要十分．
本章で用いる BDF2，AB2，EXT2 はいずれも $q=2$ で零安定であり，
初期誤差の非物理的な累積増幅を伴わない．
AB2/EXT2 の制約は零安定性ではなく，対流固有値に対する絶対安定性に現れる．

項別の時間精度設計指針：

\begin{itemize}[leftmargin=2em]
  \item \textbf{CLS 界面}：TVD-RK3 で $\Ord{\Delta t^3}$．
        $\Delta t=\min(\Delta t_{\mathrm{adv}},\Delta t_\sigma,\Delta t_{\mathrm{buoy}},\Delta t_{\mathrm{disc}})$ に律速されるため，
        TVD-RK3 の 3 段更新の時間精度がそのまま界面位置更新に反映される．ただし均質流域の全体精度は
        最小律速により流体運動方程式側 $\Ord{\Delta t^2}$ で律速される
        （詳細は \S\ref{sec:time_accuracy_table}）．
  \item \textbf{滑らかな NS ブロック}：IMEX-BDF2 + IPC + implicit-BDF2 で $\Ord{\Delta t^2}$．
        圧力--速度結合と射影閉包が均質流域での 2 次精度の上限を決める．
        変粘性界面のクロス微分項は粘性 Helmholtz DC の高次残差内で
        同時刻評価し，粘性由来の 1 次時間遅れを導入しない．
  \item \textbf{分離整合性}：本稿の fractional-step（分割ステップ）は順次適用型であり，
        一般の Lie 型演算子分割は標準的には $\Ord{\Delta t}$（Marchuk \cite{Marchuk1990}）である．
        本章で 2 次と主張するのは，IPC 補正（\S\ref{sec:ipc_derivation}）により
        \textbf{圧力投影の分離誤差} が BDF2 の時間幅 $\Delta t_{\mathrm{proj}}$ と整合する範囲である．
        CLS 先行更新・界面張力ジャンプ・浮力残差・粘性クロス項を含む項別合成の次数は，
        表~\ref{tab:time_accuracy_prod} の律速行に従う．
\end{itemize}

% -------------------------------------------------------------------------------
\subsection{演算子別剛性マップ}
\label{sec:time_operator_stiffness}
% -------------------------------------------------------------------------------

一流体式 NS 系
\[
  \rho(\psi)\,\partial_t\bu
    = -\rho(\psi)\,\mathcal{C}(\bu,\bu)
    + \mathcal{V}(\bu) - \bnabla p + \bm{f}_\sigma + \rho(\psi)\bg,
  \qquad \bnabla\cdot\bu = 0
\]
を演算子ごとに分解すると，剛性源が互いに独立であることが分かる：

\begin{center}
\renewcommand{\arraystretch}{\PaperTableArrayStretch}
\begin{tabular}{lll}
\hline
演算子 & 剛性源 & 本稿の採用構成 \\
\hline
$\mathcal{C}(\bu,\bu)$（NS 対流） & 対流 CFL $|\bu|\Delta t/h$ & IMEX-BDF2（EXT2） \\
CLS 移流 & 界面移流 CFL $|\bu|\Delta t/h$ & TVD-RK3 \\
$\mathcal{V}(\bu)$（粘性） & 粘性 CFL $\nu\Delta t/h^2$ & implicit-BDF2 + 欠陥補正（DC） \\
$\bm{f}_\sigma$（界面張力） & 界面張力 CFL $\Delta t_\sigma$ & PPE 内蔵の圧力ジャンプ補正量 \\
$\bnabla p$，$\bnabla\cdot\bu=0$ & 非圧縮拘束 & IPC + FCCD PPE \\
$\rho\bg$（浮力，上昇気泡） & 浮力波 $\Delta t \lesssim N_{\mathrm{BV}}^{-1}$ & Balanced--Force 浮力分解 \\
\hline
\end{tabular}
\end{center}

\noindent\textbf{分離整合性：}
fractional-step（分割ステップ）構成は Lie 型演算子分割（順次適用）に分類されるが，
IMEX-BDF2 の予測子 + IPC 投影 + implicit-BDF2 粘性更新は
1 物理ステップ内で同じ $\Delta t_{\mathrm{proj}}$ を共有するため，
圧力投影に由来する分離誤差を $\Ord{\Delta t^2}$ に抑える．
Strang 型対称分割~\cite{Strang1968} は 2 次精度を保証する一般原理だが，
本稿では対称化そのものではなく，IPC の増分圧力と BDF2 射影時間幅の整合により
圧力--速度結合の 2 次性を確保する．項別合成の最終次数は
表~\ref{tab:time_accuracy_prod} の律速項で決まる．

\subsubsection*{CFL 数と固有値マップ}
\label{sec:cfl_definitions}

各剛性源を Fourier モード $u_j^n = \hat{u}^n e^{ik_j h}$，$\theta\equiv kh\in[-\pi,\pi]$
で展開すると，半離散左辺 $\partial_t \hat{u} = \lambda(\theta)\hat{u}$ の
固有値 $\lambda(\theta)$ は演算子ごとに異なる軌跡を取る．
時間積分の絶対安定領域 $\mathcal{S}\subset\mathbb{C}$ に対し
$\Delta t\,\lambda(\theta)\in\mathcal{S}$（$\forall\theta$）が必要であり，
これが項別 CFL 数の幾何学的根拠となる：
\begin{equation}
  \mathrm{CFL}_{\mathrm{adv}} \equiv \frac{|\bu|\Delta t}{h},\qquad
  \mathrm{CFL}_{\nu} \equiv \frac{\nu\Delta t}{h^2},\qquad
  \mathrm{CFL}_{\sigma} \equiv \frac{\Delta t}{\Delta t_\sigma},\qquad
  \Delta t_\sigma \equiv
  \sqrt{\frac{(\rho_l+\rho_g)h_{\min}^3}{2\pi\sigma}}
  \label{eq:cfl_number}
\end{equation}
それぞれ \textbf{対流 CFL}・\textbf{粘性 CFL}・\textbf{界面張力 CFL} と呼ぶ．
ここで一様格子では $h_{\min}=h$，非一様格子では
$h_{\min}=\min(\Delta x,\Delta y)$ とし，後段の合成式でも同じ代表幅を用いる．
固有値マップは概ね次の通りである：
\begin{itemize}[leftmargin=2em]
  \item \textbf{対流 $\mathcal{C}(\bu)$}：
        非散逸中心差分の理想化極限では $\lambda_C(\theta)\in i\mathbb{R}$（純虚軸）だが，
        AB2/EXT2 は純虚軸上に実用的な絶対安定区間を持たない．
        本稿の対流安定性は UCCD6/FCCD の風上・面フラックス合成で生じる
        離散固有値の負実部を含む複素スペクトルに対して評価し，
        $\mathrm{CFL}_{\mathrm{adv}}\le C_{\mathrm{adv}}$ を
        離散演算子込みの安定上限として用いる．
  \item \textbf{粘性 $\mathcal{V}(\bu)$}：
        $\lambda_V(\theta)\in(-\infty,0]$（負実軸）．
        A-stable な BDF2 は負実軸全域を安定領域に含むため
        $\mathrm{CFL}_{\nu}$ 制約は消滅する（§\ref{sec:time_viscous}）．
  \item \textbf{界面張力 $\bm{f}_\sigma$}：
        界面波分散関係から $\lambda_\sigma\in i\mathbb{R}$（純虚軸；
        波数 $k$ に対して $\omega^2 = \sigma k^3/(\rho_l+\rho_g)$）．
        陽的な体積力処理では界面張力 CFL $\Delta t_\sigma$
        （§\ref{sec:val_capillary}）が数値安定性を律速する．
        本稿は表面張力を PPE 内蔵の圧力ジャンプ補正量として
        圧力補正と同時刻で閉じるが，界面波の時間解像条件として
        $\Delta t_\sigma$ は残る．
  \item \textbf{浮力 $\rho\bg$}：
        重力波 $N_{\mathrm{BV}}$（Brunt--Väisälä）由来；
        Balanced--Force 浮力分解により残差成分のみが陽的に時間刻み制約を生む．
\end{itemize}

\subsubsection*{von Neumann 解析と安定領域分類}
\label{sec:von_neumann}

線形定数係数の半離散方程式 $\partial_t \hat{u} = \lambda \hat{u}$ に
時間積分 $\hat{u}^{n+1} = R(z)\hat{u}^n$（$z\equiv\Delta t\,\lambda$）を適用すると，
\textbf{増幅率（amplification factor）} $G(\theta) \equiv R(\Delta t\,\lambda(\theta))$ が定まる．
\textbf{von Neumann 安定条件}は
\begin{equation}
  |G(\theta)| \le 1 \quad \forall\,\theta\in[-\pi,\pi]
  \label{eq:von_neumann_stability}
\end{equation}
であり，線形定数係数では L2 安定の必要十分条件である．
$\mathcal{S}\equiv\{z\in\mathbb{C}:|R(z)|\le 1\}$ を\textbf{絶対安定領域}と呼び，
その形状でスキームを分類する：
\begin{itemize}[leftmargin=2em]
  \item \textbf{A-stable}：左半平面 $\{z:\Re z\le 0\}$ 全域 $\subset\mathcal{S}$．
        BDF1/BDF2 が該当．粘性項に十分．
  \item \textbf{L-stable}：A-stable かつ $|R(z)|\to 0$（$|z|\to\infty$）．
        BDF1 / Radau IIA 等．剛性拡散の高周波成分減衰に有効．
  \item \textbf{A($\alpha$)-stable}：$\{z:|\arg(-z)|<\alpha\}\subset\mathcal{S}$．
        BDF$k$（$k\ge 3$）は A($\alpha$)-stable のみ（$\alpha<\pi/2$）．
  \item \textbf{弱条件付き安定（離散スペクトル依存）}：
        TVD-RK3 は虚軸上に有限の安定区間を持つ一方，
        AB2/EXT2 は純虚軸対流には安定でない．
        NS 対流の EXT2 評価では，UCCD6/FCCD による負実部付き離散固有値
        $\lambda_C^h(\theta)$ に対して
        $\Delta t\,\lambda_C^h(\theta)\in\mathcal{S}$ を満たす範囲を
        $\mathrm{CFL}_{\mathrm{adv}}$ の上限とする．
\end{itemize}
本章の安定性分類は，
（i）粘性項に A-stable の implicit-BDF2，
（ii）NS 対流に離散スペクトル依存の EXT2 外挿（IMEX-BDF2），
（iii）CLS 移流に TVD-RK3（SSP），
の 3 軸で整理できる．圧力投影・界面張力・浮力は，この安定性分類とは別に
同時刻閉包と物理解像条件として後続節で扱う．証明と $\rho$--$\sigma$ 多項式の根分布は
Hairer--Wanner~\cite{HairerWanner1996} Sec.~V.1 を参照．

安定性分類と同時刻閉包を分けて述べることが，本章の組織化原則である．

% -------------------------------------------------------------------------------
\subsection{移流項：CLS 用 TVD-RK3 + NS 用 IMEX-BDF2}
\label{sec:time_level1}
% -------------------------------------------------------------------------------

CLS 界面移流と NS 対流の時間処理は別系統で扱う：
界面位置の高精度追跡には $\Ord{\Delta t^3}$ の TVD-RK3 を
（全体精度に対する寄与は最小律速で抑制；\S\ref{sec:time_accuracy_table}），
NS 対流項には圧力投影法と相性の良い $\Ord{\Delta t^2}$ の IMEX-BDF2 を採用する．
両移流処理からは対流 CFL 制約 $|\bu|\Delta t/h \lesssim 1$ が陽的制約として残る．
章全体の時間刻みは後段の合成式で，界面張力・浮力・離散スペクトル制約との最小値として決まり，
粘性 CFL は §~\ref{sec:time_viscous} の implicit-BDF2 で消滅させる．

\subsubsection*{CLS 界面移流：TVD-RK3（SSPRK3）}

本稿の採用構成では，CLS 移流方程式 \eqref{eq:cls_advection} の物理時間更新に
Shu--Osher 型 TVD-RK3 法（= SSPRK3）\cite{ShuOsher1988} を用いる：
\begin{align}
  q^{(1)} &= q^n + \Delta t \, \mathcal{L}(q^n) \notag\\
  q^{(2)} &= \frac{3}{4}q^n + \frac{1}{4}q^{(1)} + \frac{1}{4}\Delta t\,\mathcal{L}(q^{(1)}) \notag\\
  q^{n+1} &= \frac{1}{3}q^n + \frac{2}{3}q^{(2)} + \frac{2}{3}\Delta t\,\mathcal{L}(q^{(2)})
  \label{eq:tvd_rk3}
\end{align}
ここで $\mathcal{L}(q)$ は空間演算子である．
再初期化方程式 \eqref{eq:cls_reinit_iso} を
仮想時間 $\tau$ で陽的に積分する場合にも同じ SSPRK3 形式で書けるが，
Ridge--Eikonal 補助経路は仮想時間発展を用いない距離再構成
（\S\ref{sec:xi_sdf_reinit}）であり，物理時間ステップの律速には含めない．
本稿の CLS 移流では \textbf{FCCD フラックス発散}
$\mathcal{L}(\psi)=-\bnabla\cdot(\psi\bu)$
（空間演算子 $\mathcal{L}$ の具体的な構成は §~\ref{sec:advection} で詳述）を用いる．
運動量移流 $(u,v)$ の空間スキームには UCCD6（§~\ref{sec:uccd6_def}）を用いる．

\noindent\textbf{CLS 専用 $\psi^{n+1}$ 更新：}
$q \leftarrow \psi$ かつ $\bu=\bu^n$（直前ステップの \textbf{PPE 投影済み}速度；
未投影予測値 $\bu^*$ を採ると後述の §\ref{sec:cls_velocity_consistency} で示すように
PPE 右辺不整合と $\psi$ 質量変化を生む）として
\eqref{eq:tvd_rk3} を CLS に適用すると，本稿の界面位置更新は
\begin{align}
  \psi^{(1)} &= \psi^n - \Delta t\,\bnabla\cdot(\psi^n \bu^n) \notag\\
  \psi^{(2)} &= \tfrac{3}{4}\psi^n + \tfrac{1}{4}\psi^{(1)}
               - \tfrac{1}{4}\Delta t\,\bnabla\cdot(\psi^{(1)}\bu^n) \notag\\
  \psi^{n+1} &= \tfrac{1}{3}\psi^n + \tfrac{2}{3}\psi^{(2)}
               - \tfrac{2}{3}\Delta t\,\bnabla\cdot(\psi^{(2)}\bu^n)
  \label{eq:cls_tvd_rk3_psi}
\end{align}
となる（$\bnabla\cdot$ は FCCD 面フラックスで離散化）．
$\psi^{n+1}$ は本ステップ内で\textbf{最初に確定する}変数であり，
NS 予測子（§\ref{sec:time_level1} の IMEX-BDF2）以降は
$\rho^{n+1} = \rho_g + (\rho_l - \rho_g)\psi^{n+1}$ と曲率 $\kappa_{lg}^{n+1}$
を確定値として PPE の圧力ジャンプ閉包と浮力残差に流す疎結合構成となる．
全体時間精度の閉包根拠は §\ref{sec:time_accuracy_table} を参照．

\noindent\textbf{TVD（Total Variation Diminishing）の定義：}
$TV(q) \equiv \sum_i |q_{i+1} - q_i|$ が時間的に増大しない性質を指す．
TVD-RK3 は，空間離散化を前進 Euler で進めた 1 ステップが TVD となる場合に
その性質を時間方向へ継承する SSP 型 Runge--Kutta 法である．
本稿の FCCD は線形フィルタであり空間 TVD 性を持たないため，
TVD-RK3 との組合せだけで全変動減少は保証されない（注意~\ref{warn:adv_clamp}）．
値域制限（式~\eqref{eq:psi_clamp}）と組み合わせることで $\psi \in [0,1]$ を維持する．

\begin{tcolorbox}[warnbox, title={TVD-RK3 の適用範囲}]
\label{warn:tvd_rk3_scope}
本稿の採用構成で TVD-RK3 を用いるのは CLS 保存形移流である
（NS 対流には使わない）：
\begin{itemize}
  \item \textbf{CLS 移流方程式}（式~\eqref{eq:cls_advection}）：
        $\pd{\psi}{t} + \bnabla\cdot(\psi\,\bu) = 0$（保存形）
\end{itemize}
再初期化方程式（式~\eqref{eq:cls_reinit_iso}）の仮想時間 TVD-RK3 は
仮想時間方程式として書く場合の補助形式であり，
Ridge--Eikonal 補助経路の幾何再構成とは別の時間スケールに属する．
NS 対流は次小節の IMEX-BDF2（EXT2 外挿）で扱う．
\end{tcolorbox}

% -------------------------------------------------------------------------------
\subsubsection{速度--PPE 整合性：CLS 移流が用いる速度場の正しい順序}
\label{sec:cls_velocity_consistency_v7}
\phantomsection\label{sec:cls_velocity_consistency}
% -------------------------------------------------------------------------------

CLS 保存形更新（\S\ref{sec:cls_stages}）が用いる速度は
\textbf{直前の PPE 投影済み} $\bu^n$ である．
一方，運動量移流（\S\ref{sec:fccd_advection}）の EXT2 外挿は
PPE 投影済みの履歴 $\{\bu^n,\bu^{n-1}\}$ から構成し，
未投影 $\bu^*$（運動量予測子出力）は移流速度として用いない．
未投影 $\bu^*$ を使うと：
\begin{itemize}
  \item 移流で $\bu^*$ の離散発散残差が $\psi$ 質量変化を生む
  \item CLS 移流と運動量予測子で速度場が異なり，
        それぞれに対応する PPE 右辺が不整合
\end{itemize}
本稿が採用する時間ステップ順序：
\begin{equation*}
  (\text{PPE}_{n-1},\text{PPE}_{n})\to\{\bu^{n-1},\bu^{n}\}
  \to\text{CLS 更新}\,\&\,\text{運動量予測子}
  \to\bu^*_{n+1}\to\text{PPE}_{n+1}\to\bu^{n+1}.
\end{equation*}
これは時間積分順序の根幹であり，\S\ref{sec:cls_stages} の CLS 保存形更新と
\S\ref{sec:time_viscous} の粘性陰解法・\S\ref{sec:ipc_derivation} の
IPC 圧力投影が共有する大域的な因果順序である．

\subsubsection*{NS 対流：IMEX-BDF2（EXT2 外挿）}

水--空気級の密度比・非一様格子・長時間実験では NS 対流項に
\textbf{IMEX-BDF2 法（EXT2 外挿）} を採用する．対流項は
\[
  \mathcal{C}^{\mathrm{EXT2}}
  = 2\mathcal{C}(\bu^n,\bu^n)
    - \mathcal{C}(\bu^{n-1},\bu^{n-1})
\]
で陽的外挿し，粘性ブロックは BDF2 の係数 $\gamma=2/3$ を用いて
\begin{equation}
  \frac{3\bu^* - 4\bu^n + \bu^{n-1}}{2\Delta t}
  = -\mathcal{C}^{\mathrm{EXT2}}
    + \mathcal{V}(\bu^*)
    + \bm{a}^{\,n+1}_{\mathrm{res}}
  \label{eq:predictor_imex_bdf2}
\end{equation}
を解く．ここで $\mathcal{C}(\bu,q) = \bu\cdot\bnabla q$（対流演算子），
$\bm{a}^{\,n+1}_{\mathrm{res}}$ は式~\eqref{eq:buoyancy_gradient_residual_split} の
浮力残差を投影と同じ面演算子で評価した加速度であり，静水圧勾配成分は
PPE の楕円型作用素に残す（Balanced--Force 浮力分解；§~\ref{sec:time_buoyancy}）．
Helmholtz 形で書けば
\[
  \bigl(I-\gamma\Delta t\,\mathcal{V}\bigr)\bu^*
  =
  \frac{4}{3}\bu^n-\frac{1}{3}\bu^{n-1}
  +\gamma\Delta t\bigl(-\mathcal{C}^{\mathrm{EXT2}}
  +\bm{a}^{\,n+1}_{\mathrm{res}}\bigr).
\]
続く投影補正も同じ射影時間幅
$\Delta t_{\mathrm{proj}}\equiv\gamma\Delta t$ を用いる：
\begin{equation}
  \bu_f^{n+1}
  = \bu_f^* - \Delta t_{\mathrm{proj}}\,A_f G_f(\delta p),\qquad
  D_f\bu_f^{n+1}=0,
  \label{eq:projection_imex_bdf2_dt}
\end{equation}
すなわち投影形 PPE は $D_f A_f G_f$，SPD 形は $-D_f A_f G_f$ として扱い，
補正子は同じ $A_fG_f$ を共有する．
界面張力は PPE 右辺のジャンプ条件として陰的に課す（§~\ref{sec:time_csf}）．

\paragraph{安定性：対流 CFL と IMEX-BDF2 安定領域}
本段落の安定性議論は，本章前半の CFL 数定義
（式~\eqref{eq:cfl_number}）と A-stable / 条件付き安定の分類に従う．
EXT2 は対流項のみを陽的に扱うため対流 CFL 条件
$\mathrm{CFL}_{\mathrm{adv}} = |\bu|\Delta t/h \le C_{\mathrm{adv}}$ が必要である．
粘性ブロックは BDF2 で陰的に扱うため，BDF2 が A-stable
（Hairer--Wanner~\cite{HairerWanner1996} Sec.~V.1）であることから
粘性 CFL $\mathrm{CFL}_{\nu}=\nu\Delta t/h^2$ は時間刻み制約から外れる．
EXT2 は AB2 と同じく $\Ord{\Delta t^2}$ の零安定多段法であり，
NS 対流項の離散固有値 $\lambda_C^h$ が
UCCD6/FCCD の風上・面フラックス合成により負実部を持つ範囲で
絶対安定領域に入る．したがって $C_{\mathrm{adv}}$ は
純虚軸の理論値ではなく，離散演算子込みの安定係数である．
したがって本段落が決めるのは移流由来の上限であり，最終的な時間刻みは
式~\eqref{eq:dt_per_term} の通り，対流・界面張力・浮力・離散スペクトル制約の
最小値として合成する．

\begin{tcolorbox}[warnbox, title={IMEX-BDF2 / EXT2 起動問題（$n=0,1$）}]
\label{warn:bdf2_startup}
EXT2 外挿および BDF2 は $n\ge 1$ でのみ適用可能である．
初回ステップ ($n=0$) は BDF1（後退 Euler）で 1 ステップ助走し
$n=1$ 以降に BDF2 + EXT2 へ移行する．BDF1 起動は次時刻の値に
$\Ord{\Delta t^2}$ の開始誤差を与えるが，BDF2 / EXT2 の零安定性
（Dahlquist の根条件）によりこの開始誤差は後続ステップで増幅されず，
全体精度 $\Ord{\Delta t^2}$ と整合する \cite{DahlquistBjorck1974}．
\end{tcolorbox}
